% This is text associated with the open textbook "Open Electromagnetics"
% See immediately below for author, date
% This work is licensed under the Creative Commons Attribution-ShareAlike 4.0 International License. (CC BY SA 4.0) 
% To view a copy of the license, visit http://creativecommons.org/licenses/by-sa/4.0/

%ORB TITLE Poor Conductors
%ORB AUTHOR 0        # S.W. Ellingson, Virginia Tech (ellingson@vt.edu)
%ORB DATE 20191115   # yyyymmdd
%ORB ABSTRACT

\label{m0156_Poor_Conductors} % <-- Must match name of this file and module directory name.  Don't mess with this!
\index{conductor!poor|(}

A \emph{poor conductor} is a material for which conductivity is low, yet sufficient to exhibit significant loss.
To be clear, the loss we refer to here is the conversion of the electric field to current through Ohm's law.

The threshold of significance depends on the application. 
For example, the dielectric spacer separating the conductors in a coaxial cable might be treated as lossless ($\sigma=0$) for short lengths at low frequencies; whereas the loss of the cable for long lengths and higher frequencies is typically significant, and must be taken into account.
In the latter case, the material is said to be a poor conductor because the loss is significant yet the material can still be treated in most other respects as an ideal dielectric.

A quantitative but approximate criterion for identification of a poor conductor can be obtained from the concept of complex permittivity $\epsilon_c$, which has the form:
%
\begin{equation}
\epsilon_c = \epsilon' - j\epsilon''
\end{equation}
%
Recall that $\epsilon''$ quantifies loss, whereas $\epsilon'$ exists independently of loss.
In fact, $\epsilon_c=\epsilon'=\epsilon$ for a perfectly lossless material.
Therefore, we may quantify the lossiness of a material using the ratio $\epsilon''/\epsilon'$, which is sometimes referred to as \emph{loss tangent}
\index{loss tangent}
(see Section~\ref{m0132_Loss_Tangent}).
Using this quantity, we define a poor conductor as a material for which $\epsilon''$ is very small relative to $\epsilon'$.
Thus,
%
\begin{equation}
\frac{\epsilon''}{\epsilon'} \ll 1 ~~~ \mbox{(poor conductor)}
\label{m0156_eDef}
\end{equation}
%
%%%%%%%%%%%%%%%%%%%%%%%%%%%%%%%%%%%%%%%%%%%%%%%
%ORB HIGHLIGHT
\begin{mdframed}[backgroundcolor=yellow!20,nobreak=true] 
A poor conductor is a material having loss tangent much less than 1, such that it behaves in most respects as an ideal dielectric except that ohmic loss may not be negligible. 
\end{mdframed}
%%%%%%%%%%%%%%%%%%%%%%%%%%%%%%%%%%%%%%%%%%%%%%%

An example of a poor conductor commonly encountered in electrical engineering includes the popular printed circuit board substrate material FR4
\index{FR4}
(fiberglass epoxy), which has $\epsilon''/\epsilon' \sim 0.008$ over the frequency range it is most commonly used.
Another example, already mentioned, is the dielectric spacer material (for example, polyethylene) 
\index{polyethylene}
typically used in coaxial cables.
The loss of these materials may or may not be significant, depending on the particulars of the application.

The imprecise definition of Equation~\ref{m0156_eDef} is sufficient to derive some characteristics exhibited by all poor conductors.
To do so, first recall that the propagation constant $\gamma$ is given in general as follows:
%
\begin{equation}
\gamma^2 = -\omega^2\mu\epsilon_c
\end{equation}
%
Therefore:
%
\begin{equation}
\gamma = \sqrt{-\omega^2\mu\epsilon_c}
\end{equation}
%
In general a number has two square roots, so some caution is required here. 
In this case, we may proceed as follows:
%\footnote{A recommended exercise for the student is to repeat this derivation with the other root, and observe that an absurd result is obtained.}
%
\begin{flalign}
\gamma &= j\omega\sqrt{\mu}\sqrt{\epsilon'-j\epsilon''} \nonumber \\
       &= j\omega\sqrt{\mu\epsilon'}\sqrt{1-j\frac{\epsilon''}{\epsilon'}} \label{m0156_eGE}
\end{flalign}
%
The requirement that $\epsilon''/\epsilon' \ll 1$ for a poor conductor allows this expression to be ``linearized.''
For this, we invoke the \emph{binomial series} 
\index{binomial series}
representation:
%
\begin{equation}
\left(1+x\right)^n = 1 + nx + \frac{n(n-1)}{2!} x^2 + ...
\end{equation}
%
where $x$ and $n$ are, for our purposes, any constants; and ``...'' indicates the remaining terms in this infinite series, with each term containing the factor $x^n$ with $n>2$.
If $x\ll 1$, then all terms containing $x^n$ with $n\ge 2$ will be very small relative to the first two terms of the series.
Thus,
%
\begin{equation}
\left(1+x\right)^n \approx 1 + nx ~~~\mbox{for $x\ll 1$}
\end{equation}
%
Applying this to the present problem:
%
\begin{equation}
\left(1-j\frac{\epsilon''}{\epsilon'}\right)^{1/2} \approx 1 - j\frac{\epsilon''}{2\epsilon'} 
\end{equation}
%
where we have used $n=1/2$ and $x=-j\epsilon''/\epsilon'$.
Applying this approximation to Equation~\ref{m0156_eGE}, we obtain:
%
\begin{flalign}
\gamma &\approx j\omega\sqrt{\mu\epsilon'}\left(1 - j\frac{\epsilon''}{2\epsilon'}\right) \nonumber \\
       &\approx j\omega\sqrt{\mu\epsilon'} + \omega\sqrt{\mu\epsilon'}\frac{\epsilon''}{2\epsilon'} 
\end{flalign}
%
At this point, we are able to identify an expression for the phase propagation constant:
%
\begin{equation}
\boxed{
\beta \triangleq \mbox{Im}\left\{\gamma\right\} \approx \omega\sqrt{\mu\epsilon'} ~~~\mbox{(poor conductor)}
}
\end{equation}
%
Remarkably, we find that $\beta$ for a poor conductor is approximately equal to $\beta$ for an ideal dielectric.

For the attenuation constant, we find
%
\begin{equation}
\boxed{
\alpha \triangleq \mbox{Re}\left\{\gamma\right\} \approx \omega\sqrt{\mu\epsilon'}\frac{\epsilon''}{2\epsilon'} ~~~\mbox{(poor conductor)}
}
\label{m0156_ealpha1}
\end{equation}
%
Alternatively, this expression may be written in the following form:
%
\begin{equation}
\alpha \approx \frac{1}{2}\beta\frac{\epsilon''}{\epsilon'} ~~~\mbox{(poor conductor)}
\end{equation}
%
Presuming that $\epsilon_c$ is determined entirely by ohmic loss,
%and is not accounting for delayed polarization response (as described in Section~\ref{m0134_Complex_Permittivity}), 
then
%
\begin{equation}
\frac{\epsilon''}{\epsilon'} = \frac{\sigma}{\omega\epsilon}
\end{equation}
%
Under this condition, Equation~\ref{m0156_ealpha1} may be rewritten:
%
\begin{equation}
\alpha \approx \omega\sqrt{\mu\epsilon'}\frac{\sigma}{2\omega\epsilon} ~~~\mbox{(poor conductor)}
\end{equation}
%
Since $\epsilon'=\epsilon$ under these assumptions, the expression simplifies to 
%
\begin{equation}
\alpha \approx \frac{\sigma}{2}\sqrt{\frac{\mu}{\epsilon}} = \frac{1}{2}\sigma\eta  ~~~\mbox{(poor conductor)}
\end{equation}
%
where $\eta\triangleq\sqrt{\mu/\epsilon'}$ is the wave impedance 
\index{wave impedance}
presuming lossless material.
This result is remarkable for two reasons:  
First, factors of $\omega$ have been eliminated, so there is no dependence on frequency separate from the frequency dependence of the constitutive parameters $\sigma$, $\mu$, and $\epsilon$.  These parameters vary slowly with frequency, so the value of $\alpha$ for a poor conductor also varies slowly with frequency.
Second, we see $\alpha$ is proportional to $\sigma$ and $\eta$.  This makes it quite easy to anticipate how the attenuation constant is affected by changes in conductivity and wave impedance in poor conductors.

Finally, what is the wave impedance 
\index{wave impedance}
in a poor conductor?  In contrast to $\eta$, $\eta_c$ is potentially complex-valued and may depend on $\sigma$.
First, recall:
%
\begin{equation}
\eta_c = \sqrt{\frac{\mu}{\epsilon'}} \cdot \left[ 1-j\frac{\epsilon''}{\epsilon'} \right]^{-1/2}
\end{equation}
%
Applying the same approximation applied to $\gamma$ earlier, this may be written
%
\begin{equation}
\eta_c \approx \sqrt{\frac{\mu}{\epsilon'}} \cdot \left[ 1-j\frac{\epsilon''}{2\epsilon'} \right]~~~\mbox{(poor conductor)}
\end{equation}
%
We see that for a poor conductor, Re$\left\{\eta_c\right\}\approx\eta$ and that Im$\left\{\eta_c\right\}\ll\mbox{Re}\left\{\eta_c\right\}$.
The usual approximation in this case is simply
%
\begin{equation}
\boxed{
\eta_c \approx \eta ~~~\mbox{(poor conductor)}
}
\end{equation}

%%%%%%%%%%%%%%%%%%%%%

{\bf Additional Reading:}
\begin{itemize}
\item ``\href{https://en.wikipedia.org/wiki/Binomial_series}{Binomial series}'' on Wikipedia.
\end{itemize}

\index{conductor!poor|)}


